% from ega4-11-fr.tex
\oldpage[IV-3]{165}
  remplacer tous les $s_j$ ($1\leq j\leq n$) par $s_j^m$ sans changer les
  propriétés des $\mathfrak J_h$ et des $A'_h$, et supposer par la suite que
  $m=1$. Alors $z'$, étant annulé par $\mathfrak J_{i-1}$, s'identifie à un
  élément de
  $\operatorname{Hom}_{A'}(A'/\mathfrak J_{i-1}A',M')$, et comme
  $\mathfrak J_{i-1}A'$ est un idéal de type fini de $A'$, ce module
  d'homomorphismes s'identifie lui-même à
  $\operatorname{Hom}_A(A/\mathfrak J_{i-1},M)\otimes_A A'$
  $(\mathbf 0_{\mathrm I},6.2.2)$. Soit
  \[
  v_\lambda=\operatorname{Hom}(1,u_\lambda):
  \operatorname{Hom}_A(A/\mathfrak J_{i-1},M)\longrightarrow
  \operatorname{Hom}_A(A/\mathfrak J_{i-1},N_\lambda)
  \]
  l'homomorphisme déduit de $u_\lambda$~; la famille $(v_\lambda)$ est elle
  aussi \emph{séparante}. En effet, pour tout $t\in A$, on a
  \[
  (\operatorname{Hom}_A(A/\mathfrak J_{i-1},M))_t
  =\operatorname{Hom}_{A_t}(A_t/(\mathfrak J_{i-1})_t,M_t)
  \]
  et de même en remplaçant $M$ par $N_\lambda$, puisque l'idéal
  $\mathfrak J_{i-1}$ est de type fini $(\mathbf 0_{\mathrm I},1.3.5)$~;
  comme par hypothèse l'intersection des noyaux des $(u_\lambda)_t$ est nulle,
  il en est de même de l'intersection des noyaux des
  $(v_\lambda)_t=\operatorname{Hom}(1,(u_\lambda)_t)$, d'où notre assertion
  (11.9.4). Remplaçant $A$ par $A/\mathfrak J_{i-1}$, $M$ par
  $\operatorname{Hom}_A(A/\mathfrak J_{i-1},M)$, $N_\lambda$ par
  $\operatorname{Hom}_A(A/\mathfrak J_{i-1},N_\lambda)$ (qui sont des
  $(A/\mathfrak J_{i-1})$-modules), $u_\lambda$ par $v_\lambda$ et enfin $A'$
  par $A'/\mathfrak J_{i-1}A'$, on voit qu'on peut se ramener au cas où, dans
  la situation initiale, l'élément $s=s_i\in A$ est tel que $A'_s$ soit un
  $A_s$-module libre. Or la famille des
  $(u_\lambda)_s:M_s\to(N_\lambda)_s$ est séparante par hypothèse~; comme on a
  $(u'_\lambda)_s(z'/1)=0$, il résulte de (11.9.7) que l'on a $z'/1=0$ dans
  $M'_s$~; mais cela signifie qu'il existe un entier $r$ tel que $s^rz'=0$
  dans $M'$, ce qui achève de prouver $(*_i)$ par récurrence.
\end{enumerate}
\end{proof}

\begin{remark}[11.9.11]
\label{IV.11.9.11-fr}
Bornons-nous pour simplifier au cas où $Z_\lambda=X$ pour tout $\lambda$. Il
faut noter alors que si la famille des homomorphismes
$u_\lambda:\mathcal F\to\mathcal G_\lambda$ est séparante, \emph{il ne
s'ensuit pas nécessairement} que, pour tout $x\in X$, l'intersection des
noyaux des homomorphismes
$(u_\lambda)_x:\mathcal F_x\to(\mathcal G_\lambda)_x$ soit réduite à 0. Par
exemple, soient $X$ un préschéma de Jacobson localement noethérien, de
dimension $\geq1$, et $\mathcal F$ un $\mathcal O_X$-Module cohérent~; pour
tout point fermé $x$ de $X$, et tout entier $n\geq0$,
$\mathfrak m_x^n\mathcal F$ est un $\mathcal O_X$-Module cohérent de support
contenu dans $\{x\}$. La famille des homomorphismes canoniques
$\mathcal F\to\mathcal F/\mathfrak m_x^n\mathcal F$ (où $x$ parcourt
l'ensemble $X_0$ des points fermés de $X$ et $n$ l'ensemble des entiers
$\geq0$) est \emph{séparante}~: en effet, si $t$ est une section de
$\mathcal F$ au-dessus d'un ouvert $U$ de $X$ dont les images dans les
$\Gamma(U,\mathcal F/\mathfrak m_x^n\mathcal F)$ sont toutes nulles, il en
résulte aussitôt que pour tout point fermé $x\in U$, on a $t_x=0$~; comme
l'ensemble des points fermés contenus dans $U$ est très dense dans $U$, cela
entraîne bien $t=0$ (10.2.1). Cependant, si l'on prend
$\mathcal F=\mathcal O_X$, et si $y\in X$ est un point \emph{non fermé} de
$X$, on a $(\mathcal O_X/\mathfrak m_x^n\mathcal O_X)_y=0$ pour tout point
fermé $x$ de $X$, mais $\mathcal O_{X,y}\neq0$, et l'intersection des noyaux
des homomorphismes
$\mathcal O_{X,y}\to(\mathcal O_X/\mathfrak m_x^n\mathcal O_X)_y$ est égale
à $\mathcal O_{X,y}$.
\end{remark}

\begin{lemma}[11.9.12]
\label{IV.11.9.12-fr}
Les notations étant celles de (11.9.1) et (11.9.5), supposons la famille
$(u_\lambda)$ séparante~; supposons en outre que $X$ soit un $S$-préschéma, où
$S=\operatorname{Spec}(A)$ est un schéma affine, et que
$X'=X\times_S S'$, où $S'=\operatorname{Spec}(A')$, $A'$ étant une
$A$-algèbre~; supposons enfin que les $\mathcal G_\lambda$ soient $S$-plats.
Soit $(A'_\alpha)_{\alpha\in I}$ la famille filtrante des sous-$A$-algèbres de
type fini de $A'$~; pour tout $\alpha\in I$, posons
$S''_\alpha=\operatorname{Spec}(A'_\alpha)$,
$X''_\alpha=X\times_S S''_\alpha$,
$Z''_{\alpha\lambda}=Z_\lambda\times_X X''_\alpha$, et soient
$f''_{\alpha\lambda}:Z''_{\alpha\lambda}\to X''_\alpha$,
$\mathcal F''_\alpha$, $\mathcal G''_{\alpha\lambda}$ et
$u''_{\alpha\lambda}$ les morphismes, Modules et homomorphismes de Modules
déduits de $f_\lambda$, $\mathcal F$, $\mathcal G_\lambda$ et $u_\lambda$ au
moyen du changement de base $X''_\alpha\to X$. Alors, si, pour tout
$\alpha\in I$, la famille $(u''_{\alpha\lambda})_{\lambda\in L}$ est
séparante, il en est de même de $(u'_\lambda)_{\lambda\in L}$.
\end{lemma}

\begin{proof}
Il s'agit de prouver que si une section $t'$ de $\mathcal F'$ au-dessus d'un
ouvert affine $U'$ de $X'$ est telle que $u'_\lambda(t')=0$ pour tout
$\lambda\in L$, on a $t'=0$. Si $h_\alpha:X'\to X''_\alpha$ est la projection
